\subsection{Degree Theory}

We are considering here proper maps.
Note that if $F\colon X\to Y$ is proper with $Y$ locally compact Hausdorff, then $F$ is a closed map.
If $Z\subseteq X$ is closed, then let $y\notin F(Z)$ so that $y$ has a neighborhood $U$ with compact closure.
Then $F^{-1}\cl U\subseteq X$ is compact, and thus so is $Z\cap F^{-1}\cl U$.
And so $C=F(Z\cap F^{-1}\cl U)\subseteq Y$ is compact, thus closed, and $U-C$ is an open neighborhood of $y$ disjoint from
$F(Z)$.

\bthrm

    Let $M,N$ be smooth oriented manifolds without boundary of the same dimension $n$.
    Let $F\colon M\to N$ be a proper smooth map.
    There exists a unique integer $k$, called the {\emphcolor degree} of $F$, such that the following conditions hold:
    \benum
        \item For every compactly supported $n$-form $\omega\in\Omega^n_c(N)$,
        $$ \int_MF^*\omega = k\int_N\omega . $$
        \item For a regular value $q\in N$,
        $$ k = \sum_{x\in F^{-1}q}\sgn(x) . $$
        Here $\sgn(x)=+1$ if $\d F_x$ is orientation-preserving and $-1$ otherwise.
    \eenum

\ethrm

\bproof

    By \refmath[theorem]{topcohomologycases}, cohomology classes of top forms are determined entirely by their integrals.
    Thus if we let $\theta\in\Omega^n_c(N)$ with integral $1$, we must define $k=\int_mF^*\theta$.
    Then for any other $\omega\in\Omega^n_c(N)$ we have that $[\omega]=a[\theta]$ and so
    $$ \int_MF^*\omega = a\int_MF^*\theta = ak = k\int_N\omega . $$
    So we have (1), except we do not know that $k$ is an integer.
    For that we need (2).

    We have that $F^{-1}q\subseteq M$ is a properly embedded $0$-dimensional submanifold of $M$, thus finite.
    Let $x_1,\dots,x_m$ enumerate $F^{-1}q$, and for each $i$ let $x_i\in U_i$ be a neighborhood in which $F$ is a
    diffeomorphism $U_i\to W_i\ni q$.
    We may shrink the $U_i$s and assume they are disjoint.

    Define $K=M-\bigcup_iU_i$ which is closed, so that $F(K)\subseteq N$ is closed and disjoint from $q$.
    Let $W$ be the connected component of $W_1\cap\cdots\cap W_m\cap(N-F(K))$ containing $q$ and let
    $V_i=F^{-1}W\cap U_i$.
    Then we have that $F^{-1}W=\bigsqcup_iV_i$.
    Furthermore $F$ forms a diffeomorphism $V_i\to W$ so these are the connected components of $F^{-1}W$.
    In each of these $V_i$, $F$ is either orientation-preserving or reversing.

    Let $\omega\in\Omega_c^n(W)$ have integral $1$, so that $\int_MF^*\omega=k$.
    Furthermore because $F^*\omega$ is supported within $F^{-1}W$, we have
    $$ k = \int_MF^*\omega = \sum_i\int_{V_i}F^*\omega = \sum_i\pm1 , $$
    where the sign is determined by if the diffeomorphism $F\colon V_i\to W$ preserves or reverses orientation.
    So we are done.

    If $F^{-1}q=\emptyset$ then $q$ has a neighborhood $W$ contained in the open $N-F(M)$.
    If $\omega$ is compactly supported in $W$, then $\int_MF^*\omega=0$ so that $k=0$.
    \qed

\eproof

The following properties are immediate.

\bprop

    \benum
        \item The degree of a non-surjective map is zero.
        \item The degree of the composition of two proper maps is the product of their degrees:
        $$ \deg(G\circ F) = (\deg G)(\deg F) . $$
        \item The degree of a diffeomorphism is $1$ if it preserves orientation and $-1$ otherwise.
        \item The degree of two homotopic maps are equal.
    \eenum

\eprop

Thus we may define the degree of a {\it continuous\/} map to be the degree of any smooth map homotopic to it (assuming the
manifolds are compact).

